\section{Lecture 12 (February 10, 2025, est.): From a fiber functor to a coalgebra}\label{lec12}
\noindent\textit{Scribe: Guillermo Sanmarco}


Continue with the exact faithful $k$-linear functor $\omega:\mathcal C\to\mathrm{Vec}_k$ and the objects $P_X$ and algebras $A_X$ of Lecture~\ref{lec11}.

\subsection*{The local equivalence}
\begin{theorem}\label{thm:local-abelian-reconstruction}
The functor $\widetilde\omega_X$ is an equivalence
\[
 \langle X\rangle_{\oplus}\simeq A_X\text{-}\mathrm{mod}_{\mathrm{fd}},
 \qquad
 A_X\simeq\operatorname{End}(\omega|_{\langle X\rangle_{\oplus}}).
\]
\end{theorem}
\begin{proof}
Write $A=A_X$, $P=P_X$, and $V=\omega X$. Precomposition by $a\in A$ defines a morphism on $H_X=V^*\otimes X$ whose image under $\omega$ sends $b$ to $ba$. It preserves $P$, because $\omega P=A$ is closed under multiplication and $\omega$ reflects subobject inclusions. These maps make $P$ a right $A$-module object in $\mathcal C$. Its image under $\omega$ is the right regular module. Its natural left $A$-action under $\widetilde\omega_X$ is left multiplication, so the two actions commute.

For a finite-dimensional left $A$-module $M$, form the object
\[
 Q(M)=P\otimes_A M
 =\operatorname{coker}\bigl(P\otimes A\otimes M
 \xrightarrow{\ r\otimes\mathrm{id}-\mathrm{id}\otimes\mu_M\ }
 P\otimes M\bigr).
\]
The products with vector spaces are the external products of Lecture~\ref{lec11}. This cokernel belongs to $\langle X\rangle_{\oplus}$, and exactness gives, naturally and as left $A$-modules,
\[
 \widetilde\omega_XQ(M)\simeq A\otimes_A M\simeq M.
\]

To identify the other composite, start with the evaluation morphism
\[
 P\otimes\omega X\longrightarrow X,
\]
obtained by restricting evaluation from $H_X\otimes V$. For $X^{\oplus n}$ take its direct sum. If $Y\subseteq X^{\oplus n}$, its restriction to $P\otimes\omega Y$ factors through $Y$: after applying $\omega$, this says precisely that $A$ preserves $\omega Y$. Exact faithfulness detects the required zero composite to the quotient. Passing to quotients defines evaluation for every subquotient $Y\in\langle X\rangle_{\oplus}$. The graph argument from the previous lecture makes these evaluations independent of the presentation and natural in $Y$.

They are balanced over $A$, so they induce
\[
 Q(\widetilde\omega_XY)=P\otimes_A\omega Y\longrightarrow Y.
\]
Applying $\omega$ gives the isomorphism $A\otimes_A\omega Y\to\omega Y$. An exact faithful functor reflects isomorphisms, by applying it to kernel and cokernel. Thus this evaluation is an isomorphism. The two natural isomorphisms exhibit $Q$ as a quasi-inverse, proving the equivalence.

Finally, $A$ acts naturally on all these vector spaces, giving a map to the endomorphism algebra of the restricted functor. Conversely, a natural endomorphism is determined by its component on $X$, since all objects are subquotients of finite sums of $X$. Naturality forces that component to preserve every $\omega Y\subseteq\omega(X^{\oplus n})$, so it lies in $A$. These constructions are inverse algebra maps.
\end{proof}

\subsection*{Passing from local algebras to one coalgebra}
Set $B_X=A_X^*$. The preceding theorem and algebra--coalgebra duality give
\[
 \langle X\rangle_{\oplus}\simeq\operatorname{Comod}_{\mathrm{fd}}(B_X).
\]
If $\langle X\rangle_{\oplus}\subseteq\langle Z\rangle_{\oplus}$, restriction of natural endomorphisms gives $A_Z\to A_X$. This map is surjective. Indeed, under the local equivalence for $Z$, put $M=\omega X$ and $D=A_Z/\operatorname{Ann}(M)$. The $D$-module $M$ is faithful. Choosing a basis $m_1,\ldots,m_d$ embeds the regular module by
\[
 D\longrightarrow M^{\oplus d},\qquad a\longmapsto(am_1,\ldots,am_d).
\]
Every finite-dimensional $D$-module is a quotient of a finite direct sum of copies of $D$, hence a subquotient of copies of $M$. Thus the subcategory generated by $M$ is exactly $D\text{-}\mathrm{mod}_{\mathrm{fd}}$, and $A_X\simeq D$. Dualizing gives an injection $B_X\hookrightarrow B_Z$.

Direct sums of objects make these subcategories a directed system. Define
\[
 B=\varinjlim_X B_X=\bigcup_X B_X
\]
with the compatible coalgebra structures. Every object and every finite collection of morphisms of $\mathcal C$ belongs to some $\langle X\rangle_{\oplus}$. The local equivalences therefore assemble into a fully faithful functor $\mathcal C\to\operatorname{Comod}_{\mathrm{fd}}(B)$. It is essentially surjective: the finitely many matrix coefficients of a finite-dimensional $B$-comodule lie in some $B_X$, so its coaction already factors through that subcoalgebra. We have proved the abelian reconstruction theorem
\begin{equation}\label{eq:abelian-reconstruction}
 \mathcal C\simeq\operatorname{Comod}_{\mathrm{fd}}(B),
\end{equation}
compatibly with $\omega$.

\subsection*{Natural endomorphisms and matrix coefficients}
Let $A$ now be any finite-dimensional algebra and let $U:A\text{-}\mathrm{mod}_{\mathrm{fd}}\to\mathrm{Vec}_k$ forget the action. For a commutative algebra $R$, put $U_R=U(-)\otimes R$. There is an algebra isomorphism
\[
 A\otimes R\simeq\operatorname{End}(U_R).
\]
One direction lets $A\otimes R$ act. For the other, send a natural endomorphism $\lambda$ to $\lambda_A(1\otimes1)$. For $m\in M$, the $A$-linear map $A\to M$, $a\mapsto am$, gives
\[
 \lambda_M(m\otimes1)=\lambda_A(1\otimes1)\cdot m.
\]
This proves that the two constructions are inverse. Notice that the components of $\lambda$ need only be $R$-linear, not $A$-linear.

Applying this locally and passing to inverse limits yields
\[
 \operatorname{End}(\omega_R)
 \simeq\varprojlim_X(A_X\otimes R)
 \simeq\operatorname{Hom}_k(B,R).
\]
The multiplication on the right is convolution using the coproduct of $B$.

It is useful to describe $B$ without the directed system. For $v\in\omega X$ and $\phi\in(\omega X)^*$, write $[\phi\otimes v]_X$ for a matrix coefficient. Then
\[
 B\simeq
 \left(\bigoplus_{X}(\omega X)^*\otimes\omega X\right)\big/\mathcal R,
\]
where a small skeleton is used and $\mathcal R$ is spanned by
\[
 [\phi\otimes\omega(f)v]_Y-[\phi\circ\omega(f)\otimes v]_X,
 \qquad f:X\to Y.
\]
Indeed, a linear map from this quotient to any vector space encodes precisely a natural family of linear maps on the spaces $\omega X$ with values in their tensor products with that vector space. On each local subcategory the preceding calculation identifies the quotient with $A_X^*$; taking the directed union gives $B$. If $e_i$ is a basis of $\omega X$ with dual basis $e_i^*$, the coalgebra structure and universal coaction are
\[
\begin{split}
 \Delta[\phi\otimes v]_X
 &=\sum_i[\phi\otimes e_i]_X\otimes[e_i^*\otimes v]_X,\\
 \varepsilon[\phi\otimes v]_X&=\phi(v),\\
 \rho_X(v)&=\sum_i e_i\otimes[e_i^*\otimes v]_X.
\end{split}
\]
The expressions are independent of the basis by the canonical tensor representing the identity.

In particular, applying this construction to the forgetful functor on finite-dimensional comodules of a coalgebra $C$ recovers $C$: its finite-dimensional coefficient subcoalgebras exhaust it. A coalgebra map $C\to C'$ gives a functor on comodules by composing coactions. Conversely, a functor compatible with the forgetful functors sends the coefficient $[\phi\otimes v]_X$ to $[\phi\otimes v]_{F(X)}$; naturality makes this well-defined and the formulas above show it is a coalgebra map. This is the functorial form of abelian reconstruction used next.
